% !TeX root = ../测度论笔记.tex
% 由主文档通过 \input 载入；请编译主文档。
\setcounter{chapter}{0} % 保留原书章号。
\chapter{绪论：从长度、积分到随机性}
\label{chap:mt-introduction}

\section{为什么需要测度论}
区间的长度、平面图形的面积、集合中元素的个数，以及随机事件的概率，看似属于不同问题，却都有共同要求：空集的大小为零，大小不能为负，互不相交部分合起来的大小应当相加。\emph{测度论（measure theory）}从这些要求出发，研究如何给集合赋予大小，以及如何借助这种大小定义积分、处理函数列的极限。它既是现代概率论的基础，也是研究函数空间、弱导数和偏微分方程的重要工具。

有限次相加并不足以处理分析中的极限。例如，有理数集可以写为可数个单点的并；单点长度为零，应当推出有理数集长度为零。但实数区间是不可数个单点的并，区间长度却不为零。这里需要准确区分有限、可数与不可数，而不能把“每一部分为零”无条件推广为“总量为零”。

积分也有类似困难。函数列即使处处收敛，其积分仍可能不收敛到极限函数的积分。测度论将“如何测量集合”和“何时可以交换积分与极限”联系起来，使这些问题能够在统一条件下讨论。

\begin{example}[移动到小区间上的质量]
在 $(0,1)$ 上令 $f_n(x)=n\mathbf1_{(0,1/n)}(x)$，其中 $\mathbf1_A$ 在 $A$ 上取 $1$，在其余位置取 $0$。对每个固定的 $x>0$，充分大的 $n$ 满足 $1/n<x$，所以 $f_n(x)\to0$。然而，每个函数的积分都是高度与底边长度的乘积：
\[
 \int_0^1 f_n(x)\,\mathrm dx=1.
\]
逐点观察能看到每个位置最终归零，却不能阻止质量集中在越来越窄的区间中。第\ref{chap:mt-convergence}章将说明，交换极限与积分需要额外条件。
\end{example}

\section{定位、先修知识与章节路线}
本书面向学过数学分析基础的读者。先修知识包括集合运算、实数的上确界性质、数列与级数、连续函数、黎曼积分和函数列的逐点收敛。涉及开集时只需要欧氏空间中的初等拓扑；抽象拓扑不是起点。高等代数有助于理解函数空间，但不是定义测度和积分的必要前提。

本书以实值函数为主，非负积分允许取 $+\infty$；对一般有符号函数，始终先检查可积性。核心结果给出证明，包括测度扩张、乘积测度和 Hahn 分解；积分与微分专题中引用的较深结果则明确标出证明范围，便于区分已经证明的结论与进一步阅读的内容。

\begin{longtable}{L{0.10\textwidth}L{0.30\textwidth}L{0.48\textwidth}}
\toprule
章节 & 主题 & 需要解决的问题\\
\midrule
\endhead
\ref{chap:mt-sets} & 集合与 $\sigma$ 代数 & 哪些集合可以作为测量对象？\\
\ref{chap:mt-measures} & 测度空间与零测集 & 可列可加性怎样控制集合极限？\\
\ref{chap:mt-construction} & 外测度与测度构造 & 区间长度怎样扩张为勒贝格测度？\\
\ref{chap:mt-functions} & 可测函数与简单逼近 & 哪些函数可以进入积分理论？\\
\ref{chap:mt-integral} & 勒贝格积分 & 如何统一求和与连续积分？\\
\ref{chap:mt-convergence} & 积分收敛定理 & 极限、求和与积分何时可以交换？\\
\ref{chap:mt-modes} & 函数列的收敛方式 & 几乎处处、依测度与均值收敛有何区别？\\
\ref{chap:mt-product} & 乘积测度与迭代积分 & 多变量积分何时可以改变次序？\\
\ref{chap:mt-lp} & $L^p$ 空间 & 如何用积分衡量函数距离并保证完备性？\\
\ref{chap:mt-rn} & 有符号测度与密度 & 一个测度何时能表示成对另一个测度的积分？\\
\ref{chap:mt-calculus} & 积分、变差与微分 & 原函数与密度为何需要绝对连续性？\\
\ref{chap:mt-probability} & 与概率论的衔接 & 分布、期望与条件期望如何统一定义？\\
\bottomrule
\end{longtable}

首次学习可先到第\ref{chap:mt-convergence}章，掌握测度、可测函数、积分以及三个主要收敛定理，再结合概率论中的期望与泛函分析中的函数空间继续阅读。第\ref{chap:mt-rn}章为一般条件期望提供存在性依据。概率论中的大数定律、中心极限定理和随机过程，以及泛函分析中的算子与对偶理论，应在相应课程中进一步学习。

\section{记号、证明习惯与参考资料}
$\mathbb N=\{1,2,\ldots\}$，$\mathcal P(S)$ 表示集合 $S$ 的幂集，即全体子集组成的集合。一般测度空间写作 $(S,\mathcal A,\mu)$；补集均相对于当前底空间 $S$。$\mathcal B(\mathbb R)$ 表示实数轴的 Borel $\sigma$ 代数，$\lambda$ 表示勒贝格测度。概率空间写作 $(\Omega,\mathcal F,\mathbb P)$，随机变量通常用 $X,Y$，取值用 $x,y$。

对集合列写 $A_n\uparrow A$，表示 $A_n\subseteq A_{n+1}$ 且 $\bigcup_nA_n=A$；$A_n\downarrow A$ 类似地表示递减并以交集为 $A$。对函数列，$f_n\uparrow f$ 表示每点单调递增且极限为 $f$。除非另有说明，积分所涉及的函数均可测。

本书将“定义是否成立”“结论是否成立”和“某个计算能否实施”分开处理。遇到极限与积分交换，先列出使用的收敛定理及其假设；遇到密度，先说明相对于哪个测度；遇到几乎处处等式，先说明例外集合相对于哪个测度为零。进一步阅读可使用 MIT 的测度与积分公开讲义\cite{mit-measure}和 Teschl 的实分析教材主页及其章节索引\cite{teschl-ra}。

\section{从初等积分到一般积分的思想变化}
在数学分析中，黎曼积分先把自变量所在区间分成许多小区间，再用每段上的代表函数值与区间长度相乘，最后研究分割变细时的极限。这个构造直接保留了面积直觉，对连续函数、分段连续函数和许多显式计算十分有效。然而，在研究函数列时，极限函数可能具有复杂的不连续点；在研究随机变量时，底空间甚至没有自然的区间结构。因此，一般积分需要依赖集合的大小，而不能始终依赖把自变量轴分成小区间。

勒贝格积分首先处理示性函数：若函数在集合 $A$ 上等于一、在其余位置等于零，其积分就应等于 $A$ 的测度。对只取有限个值的函数，再按每个值对应集合的大小加权求和。一般非负函数通过越来越精细的简单函数从下方逼近。这种构造常被概括为“按函数值分类”，但这一说法只是动机；严格定义仍需要可测性、下方逼近和上确界，不能单凭图形理解替代。

两种积分之间不是互相排斥的关系。对有界闭区间上的黎曼可积函数，勒贝格积分给出相同结果，因此原有的多项式、指数函数和三角函数积分计算仍然有效。新理论的主要价值在于扩大可处理函数的范围，并为积分与极限的交换提供稳定的条件。一个积分可以很难写出初等原函数，却仍有良好的存在性、连续性和收敛性质；这些性质通常比显式表达式更重要。

还应区分“积分理论”和“积分技巧”。换元、分部积分、对称性与递推关系用于求出具体数值；测度论首先回答被积函数是否可测、积分是否存在、能否分拆以及能否交换次序。计算中如果省略这些问题，可能得到形式上相似而实际不相等的两个迭代积分。后文的每种交换定理都应当理解为对合法运算范围的说明。

\section{需要复习的分析语言}
本科生初学时的困难常来自量词，而不完全来自新概念。函数列逐点收敛要求：对每个固定自变量和每个误差容许量，存在足够大的序号，使其后各项误差足够小。这个序号可以依赖自变量。一致收敛则要求同一个序号适用于整个讨论集合。几乎处处收敛先允许排除一个零测集，再在余下每个位置使用逐点收敛的定义。三种表述中，“先选什么，后选什么”决定了不同强度的结论。

上确界和最大值也必须区别。上确界只要求是最小上界，不要求属于原集合。定义非负积分时，下方简单函数的积分可能没有最大值，但它们的上确界仍然存在于 $[0,\infty]$。例如，用有限小数从下方逼近一个无理数，任何有限阶段都没有达到目标，却可以用上确界刻画目标。积分构造中的逼近采用相同的逻辑，不需要假设某个“最佳简单函数”存在。

级数在测度论中承担了控制可数误差的作用。给定总误差 $\varepsilon>0$，把第 $n$ 步误差安排为小于 $\varepsilon2^{-n}$，便能保证全部误差之和不超过 $\varepsilon$。外测度的次可加性、Egorov 定理、几乎处处收敛子列的构造以及 $L^p$ 完备性证明都会反复使用这一安排。理解这种误差预算，比记住某一次证明中恰好出现的常数更有帮助。

紧性则把局部或无限的信息转化为有限控制。在证明闭区间的外测度等于长度时，可数开覆盖先通过紧性缩减为有限子覆盖，之后才能使用有限个区间的几何比较。紧性并不意味着集合元素有限，而是说某类无限覆盖具有有限的控制子族。学习时应同时保留这两个层次，避免把分析中的紧集误解为离散的小集合。

\section{贯穿全书的模型与证明方法}
本书反复比较三个空间。第一个是带计数测度的 $\mathbb N$，此时积分就是级数，每个单点都有质量一；第二个是带长度测度的有限区间，此时存在任意小的正测度集合，但总质量有限；第三个是整个实数轴，它允许质量向无穷远移动。许多结论在这三个空间中的差别，分别暴露了原子、有限性和逃向无穷远的影响。学习一个新定理时，可以主动将它放回这三个模型检查条件。

对新的测度恒等式，通常先在容易计算的集合族上验证，再研究使该恒等式成立的集合族是否对适当运算封闭，从而推广到整个 $\sigma$ 代数。对新的积分恒等式，则常按示性函数、简单函数、非负可测函数、可积有符号函数的顺序推广。前一种方法解决“哪些集合上成立”，后一种方法解决“哪些函数上成立”；二者结合构成很多抽象定理的证明骨架。

反例也应按结构分析。若删去有限测度条件，可以尝试让集合向无穷远移动；若删去可积控制，可以尝试让函数在越来越小的集合上越来越高；若把子列结论误写成整列结论，可以让误差在不同位置循环出现。这样的构造并非零散技巧，而是在检查定理究竟阻止了哪一种失效机制。

第一遍阅读宜先建立定义、典型例子和核心定理之间的联系，第二遍再完整复述扩张、完备性和密度存在性的证明。对每个定理，至少应能说明它的输入对象、必要假设、输出结论及一个具体用途；对每个重要假设，应尝试找出一个删去后失败的情形。习题既用于训练计算，也用于检验量词、集合族和例外零测集是否已经被正确理解。

